\section{制度工程落到户、仓与道路}
\label{sec:sui-state-capacity-local-burdens}

隋的再统一常在受禅和灭陈两件大事之间被讲完，然而新朝能否留下来，取决于命令如何穿过户籍、仓储、州县和道路。开皇律令为官府提供共同的刑名，江南接收又要求重新登记军民、改置州县、处理赋役；两者不是一条从“法治”到“服从”的直线。律文记录的是国家所要实行的规范，户籍和征役则是地方家庭在不同季节、不同官吏面前实际承受的事务。现存官修史书尤其难保存逃亡、拖欠、互助和谈判的日常，只在江南起事等节点让这些摩擦短暂露出。%
\footnote{开皇律、灭陈及江南反隋活动见账本 ST-583-KAIHUANG-CODE、ST-588-SOUTHERN-CAMPAIGN、ST-589-CHEN-CONQUEST、ST-590-JIANGNAN-REBELLION。可参《隋书》刑法志、高祖纪及杨素传、《北史》与《资治通鉴》；制度志和军功传不能替代普通户的连续记录。}

东都、通济渠和郡制改革把这种接收推进到更大的尺度。洛阳的宫城与仓储、连接江淮的水路、改称为郡的行政层级，都是为了让关中、关东和江南能够被同一套调度语言组织起来。运河遗址能够说明工程有其具体的河道、闸口和维护史，却不能证明某一年征发了多少劳力；地理志中的郡名能够说明名目变化，也不能使州县财政边界和地方关系同步变得整齐。%
\footnote{东都、通济渠与郡制见 ST-605-EASTERN-CAPITAL、ST-605-TONGJI-CANAL、ST-607-COMMANDERY-REFORM，依据《隋书·炀帝纪上》《隋书·食货志》《隋书·地理志》《通典》及隋唐洛阳、运河考古报告。工程年代、修缮期和使用强度必须分开。}

因此，炀帝时期的大型营建和远征不能只被解释成一个皇帝的奢侈。它们确实要求更密集的劳役、粮运和军事动员，而山东起事、黎阳政变与江都兵变也确实显示这套调度在不同环节失去回应；可是，给失败寻找单一的道德原因，会抹掉每一环中家庭、军队、仓储和边防的具体限制。隋留下的道路、法律和行政资源后来仍被唐朝使用；它留下的更深问题，则是国家能力何时成为地方可承受的秩序，何时又成为逃亡和武装聚集的压力。%
\footnote{王薄、杨玄感与江都兵变见 ST-611-WANG-BO-REBELLION、ST-613-YANG-XUANGAN、ST-618-JIANGDU-COUP。依据《隋书》炀帝纪、王薄传、杨玄感传、宇文化及传及《资治通鉴》；参与者规模和动机大多只能透过中央叙事看见。}
